\documentclass[12pt,a4paper]{article}
\usepackage[margin=2cm]{geometry}
\usepackage{amsmath,amssymb}
\usepackage[dvipsnames]{xcolor}
\usepackage{mdframed}
\usepackage{booktabs}
\usepackage{array}
\usepackage{enumitem}
\usepackage{hyperref}
\usepackage{parskip}
\usepackage{titlesec}

\titleformat{\section}{\large\bfseries}{}{0em}{}[\titlerule]
\titleformat{\subsection}{\normalsize\bfseries}{}{0em}{}

\setlist[itemize]{topsep=2pt,itemsep=1pt,parsep=0pt}
\setlist[enumerate]{topsep=2pt,itemsep=1pt,parsep=0pt}

\newmdenv[
  backgroundcolor=orange!15,linecolor=orange!70,linewidth=1.5pt,
  innerleftmargin=8pt,innerrightmargin=8pt,innertopmargin=6pt,innerbottommargin=6pt,
  skipabove=8pt,skipbelow=4pt]{qbox}

\newmdenv[
  backgroundcolor=blue!8,linecolor=blue!50,linewidth=1.5pt,
  innerleftmargin=8pt,innerrightmargin=8pt,innertopmargin=6pt,innerbottommargin=6pt,
  skipabove=4pt,skipbelow=8pt]{solbox}

\newmdenv[
  backgroundcolor=yellow!20,linecolor=orange!60,linewidth=1pt,
  innerleftmargin=8pt,innerrightmargin=8pt,innertopmargin=5pt,innerbottommargin=5pt,
  skipabove=4pt,skipbelow=8pt]{notesbox}

\hypersetup{colorlinks=true,linkcolor=blue!70!black}

\title{\textbf{CSE 719 — Distributed \& Cloud Computing}\\
\large Lecture 5: Replication Control, 2PC \& Multicast Ordering\\
\normalsize Question Analysis \& Precise Solutions (2022–2024)}
\author{Dr.\ Atiqur Rahman · University of Chittagong}
\date{Exam: 1 July 2026}

\begin{document}
\maketitle
\tableofcontents
\newpage

%----------------------------------------------------------------------
\section{Reference: Core Concepts}
%----------------------------------------------------------------------

\subsection{Replication}
\textbf{Definition:} An object has identical copies (replicas), each maintained by a separate server.

\textbf{Why replication?}
\begin{itemize}
  \item \textbf{Fault-tolerance:} $k$ replicas $\Rightarrow$ tolerate failure of any $(k-1)$ servers.
  \item \textbf{Load balancing:} read/write load spread across $k$ replicas $\Rightarrow$ load reduced by factor $k$.
  \item \textbf{Availability:} $1 - f^k$ vs $(1-f)$ with no replication
        (where $f$ = per-server failure probability).
\end{itemize}

\textbf{Availability table:}
\begin{center}
\renewcommand{\arraystretch}{1.2}
\begin{tabular}{cccc}
\toprule
$f$ & No replication & $k=3$ & $k=5$ \\
\midrule
0.1 & 90\% & 99.9\% & 99.999\% \\
0.05 & 95\% & 99.9875\% & 7 Nines \\
0.01 & 99\% & 99.9999\% & 10 Nines \\
\bottomrule
\end{tabular}
\end{center}

\textbf{Two challenges:}
\begin{enumerate}
  \item \textbf{Replication transparency:} clients should not know multiple replicas exist.
        Achieved via \emph{front-end} servers that route requests.
  \item \textbf{Replication consistency (one-copy serializability):}
        concurrent transactions on replicated objects must behave as if executed serially
        on a single logical copy.
\end{enumerate}

\subsection{Passive vs Active Replication}
\begin{center}
\renewcommand{\arraystretch}{1.25}
\begin{tabular}{lll}
\toprule
& \textbf{Passive (Primary-Backup)} & \textbf{Active} \\
\midrule
Update path & All writes $\to$ primary $\to$ backups & All FEs multicast to all replicas \\
Ordering & Primary provides total order & Requires ordered multicast \\
Failure & Run election if primary fails & Continue — no single point \\
Correctness & Primary enforces serial order & Total-order multicast + replicated state machine \\
\bottomrule
\end{tabular}
\end{center}

Active replication requires each FE to multicast to the entire replica group.
Ordering choices: FIFO, Causal, Total, or Hybrid.
For strong consistency: Total-order (or Hybrid-Total) + Replicated State Machine.
For fault handling: Virtual synchrony (reliable ordered multicast over changing membership).

\subsection{Multicast Ordering Types (for 2024 Q4c)}
\begin{center}
\renewcommand{\arraystretch}{1.3}
\begin{tabular}{lp{6cm}p{5cm}}
\toprule
\textbf{Ordering} & \textbf{Guarantee} & \textbf{Use case} \\
\midrule
\textbf{FIFO} & Messages from \emph{same sender} delivered in send order at all receivers & Best-effort; sender-local order only \\
\textbf{Causal} & If $M_1$ causally precedes $M_2$, all receivers deliver $M_1$ before $M_2$ & Collaborative editing, social feeds \\
\textbf{Total (Atomic)} & All receivers deliver \emph{all} messages in the same order & Replicated state machines; requires coordinator or token \\
\textbf{Hybrid (*-Total)} & Total ordering combined with FIFO or Causal (e.g., FIFO-Total) & Active replication with causality \\
\bottomrule
\end{tabular}
\end{center}

\textbf{Why ordering matters for overlapping groups:}
Two groups $G_1 = \{A, B\}$ and $G_2 = \{B, C\}$ both receive messages.
Without total ordering, node $B$ (in both groups) may see messages in a different order
from $A$ or $C$, causing inconsistency. Total ordering ensures all nodes — including
those in multiple overlapping groups — agree on delivery order.

\subsection{Two-Phase Commit (2PC)}
\textbf{Problem:} A transaction $T$ touches objects on servers $S_1 \ldots S_n$.
Must ensure atomicity: \textbf{all} commit or \textbf{none} commit (= Atomic Commit problem = Consensus).

\textbf{Phase 1 — Prepare:}
\begin{enumerate}
  \item Coordinator sends \texttt{PREPARE} to all participating servers.
  \item Each server writes tentative updates to \textbf{stable storage} (disk) — right before voting.
  \item Each server responds \texttt{YES} (can commit) or \texttt{NO} (cannot commit).
\end{enumerate}

\textbf{Phase 2 — Commit or Abort:}
\begin{enumerate}
  \item If coordinator receives \texttt{YES} from \emph{all} servers within timeout:
        send \texttt{COMMIT}; each server applies updates from disk to store; replies \texttt{OK}.
  \item If any \texttt{NO} vote \emph{or} timeout before all votes:
        send \texttt{ABORT}; each server discards tentative updates.
\end{enumerate}

\textbf{Key invariants:}
\begin{itemize}
  \item A server that voted \texttt{YES} \emph{cannot} commit or abort unilaterally
        — it must wait for the coordinator's decision.
  \item A server that voted \texttt{NO} can abort immediately (it knows the transaction will abort).
  \item Saves before voting $\Rightarrow$ crash-safe: updates survive a server crash and
        can be applied (or rolled back) after coordinator sends Phase 2 decision.
\end{itemize}

\textbf{Failure handling in 2PC:}
\begin{center}
\renewcommand{\arraystretch}{1.2}
\begin{tabular}{ll}
\toprule
\textbf{Failure type} & \textbf{Response} \\
\midrule
Server crash (before voting) & Server aborts on restart (no vote committed) \\
Server crash (after YES vote) & On restart: wait for coordinator or query its decision \\
Prepare message lost & Server timeouts $\Rightarrow$ votes NO $\Rightarrow$ coordinator aborts \\
YES/NO message lost & Coordinator timeouts $\Rightarrow$ aborts (pessimistic) \\
Commit/Abort message lost & Server polls coordinator repeatedly until answer received \\
Coordinator crash & Logs all decisions on disk; new coordinator (re-elected) takes over \\
\bottomrule
\end{tabular}
\end{center}

\subsection{2PC vs 1PC}
\textbf{One-phase commit shortcomings:}
\begin{itemize}
  \item Server with the object has no say in whether to commit or abort.
        If data is corrupted, the server cannot commit — but other servers already did.
  \item A server may crash before receiving the commit message, leaving its updates in memory (lost).
\end{itemize}
2PC fixes both: Phase 1 gives every server a vote; data is persisted to disk before voting.

\newpage
%======================================================================
\section{2024 — Q2a (1.5 marks, partial — consensus definition)}
%======================================================================

\begin{qbox}
\textbf{Q2(a):} Define consensus algorithms. Differentiate between logical concurrency
and physical concurrency. \hfill \textit{[1.5 marks]}
\end{qbox}

\begin{solbox}
\textbf{Consensus Algorithm.}
A consensus algorithm allows a set of distributed processes to \textbf{agree on a single value}
(or decision), even when some processes may fail or messages may be lost.
All non-faulty processes must eventually decide the same value, and the decided value
must have been proposed by at least one process.
\textit{Example:} Paxos, Raft, 2PC (for atomic commit).

\medskip
\textbf{Logical Concurrency.}
Multiple tasks (threads, coroutines) are \emph{logically active at the same time} but
may be time-multiplexed on a single CPU.
The OS scheduler switches between them; no true simultaneous execution is required.
\textit{Example:} 100 threads on a single-core CPU.

\medskip
\textbf{Physical Concurrency.}
Multiple tasks execute on \emph{truly separate hardware units} (multi-core, multi-CPU,
multi-machine) simultaneously.
\textit{Example:} MapReduce running 1000 Map tasks in parallel across 1000 machines.

\medskip
\textbf{Key difference:}
Logical concurrency is about the \emph{appearance} of simultaneous progress (interleaving);
physical concurrency is \emph{actual} simultaneous execution.
In distributed systems, the combination of both (logical scheduling + physical parallelism)
is what makes consistency control non-trivial: even logically sequential operations can
arrive at servers out of order.
\end{solbox}

\begin{notesbox}
\textbf{Exam pattern:} 2024 Q2a (1.5 marks). Very short answer expected.
Name consensus, give one-line definition, one-line example. Then 2 sentences on
logical vs physical. Do not over-explain.
\end{notesbox}

\newpage
%======================================================================
\section{2022 — Q3a, Q3b, Q3c (3+3+3 = 9 marks)}
%======================================================================

\begin{qbox}
\textbf{Q3:} Consider a scenario where a system handles ACID transactions.
Discuss the shortcoming of each of the following approaches when a crash occurs:\\
\textbf{(a)} Writing updates directly to disk for every transaction. \hfill\textit{[3 marks]}\\
\textbf{(b)} Buffering updates in memory and flushing to disk every 50 transactions. \hfill\textit{[3 marks]}\\
\textbf{(c)} Keeping data in memory and maintaining a separate log file on disk. \hfill\textit{[3 marks]}
\end{qbox}

\begin{solbox}
\textbf{(a) Direct-to-Disk Write Per Transaction.}

Each transaction writes its updates immediately and directly to disk before committing.

\textbf{Shortcomings:}
\begin{enumerate}
  \item \textbf{Extremely slow performance.}
        Disk I/O is $\sim\!100{,}000\times$ slower than memory access.
        Every write operation blocks until the disk controller confirms the write.
        For high-throughput systems this makes each transaction take milliseconds instead
        of microseconds.

  \item \textbf{No atomicity on partial writes.}
        A transaction that touches multiple objects writes them one at a time.
        If a crash occurs after some objects are written but before all are, the database
        is left in a \emph{partial state}: some updates are on disk and others are lost.
        There is no mechanism to undo the already-written objects.

  \item \textbf{No rollback capability.}
        There is no journal or undo log. If the transaction must be aborted (e.g., due to
        a conflict), the already-committed disk writes cannot be rolled back.
\end{enumerate}
\end{solbox}

\begin{solbox}
\textbf{(b) Memory Buffer — Flush to Disk Every 50 Transactions.}

Updates are held in main memory and written to disk only after every 50 transactions complete.

\textbf{Shortcomings:}
\begin{enumerate}
  \item \textbf{Durability violation (violates D in ACID).}
        If the system crashes before the batch flush, all up to 50 buffered transactions
        are lost entirely, even if they had already returned "committed" to the client.
        Clients believe their data is safe; it is not.

  \item \textbf{Large inconsistency window.}
        The gap between the last flush and the current time represents a window of lost
        work. The larger the batch size, the more data is at risk.

  \item \textbf{State after crash is inconsistent.}
        On restart, the system knows the state as of the last flush point, but that state
        may not reflect all committed transactions — creating a mismatch between what
        clients were told (committed) and what actually persists.

  \item \textbf{Unfair: slow transactions penalise fast ones.}
        Transactions 1–49 cannot be considered durable until transaction 50 triggers the
        flush — their durability window depends on the timing of the 50th transaction.
\end{enumerate}
\end{solbox}

\begin{solbox}
\textbf{(c) Memory Data + Disk Log File (Write-Ahead Log).}

All active data lives in memory; every update is first appended to a durable log file
on disk before being applied to in-memory state (write-ahead logging, WAL).

\textbf{Shortcomings:}
\begin{enumerate}
  \item \textbf{Long recovery time after crash.}
        On restart, the system must \emph{replay} all log entries from the last checkpoint
        to reconstruct in-memory state. For large or busy systems, this log replay can take
        minutes or longer, causing extended downtime.

  \item \textbf{Log file grows unboundedly.}
        Without periodic compaction (checkpointing), the log grows without limit.
        Checkpointing itself is expensive: it requires flushing in-memory state to disk
        and may block operations.

  \item \textbf{Double-write overhead.}
        Every write incurs two disk operations: one to the log and one to eventually
        persist the data itself. This doubles disk I/O under write-heavy workloads.

  \item \textbf{Single log file is a bottleneck / single point of failure.}
        If the log disk fails, all in-memory state since the last checkpoint is
        unrecoverable. The log must reside on reliable (possibly replicated) storage,
        increasing complexity and cost.
\end{enumerate}

\textbf{Why WAL is still the best of the three:}
Despite its shortcomings, WAL (approach c) is superior to (a) and (b) because it
provides \emph{durability} (log is on disk) and \emph{atomicity} (can replay or undo
from log) while keeping active data operations fast (in-memory). The shortcomings of
(a) and (b) are more fundamental.
\end{solbox}

\begin{notesbox}
\textbf{Exam pattern:} 2022 Q3a–c (3 marks each). Each part expects $\sim$2–3 distinct
shortcomings with brief explanations. The examiner is testing understanding of the
ACID Durability property and why naive persistence approaches fail.
Use terms: \textbf{ACID, Durability, atomicity, rollback, write-ahead log, checkpoint,
inconsistency window}.
\end{notesbox}

\newpage
%======================================================================
\section{2024 — Q2c (3 marks)}
%======================================================================

\begin{qbox}
\textbf{Q2(c):} What is the atomic commit protocol? Explain the two-phase commit protocol
for nested transactions. \hfill \textit{[3 marks]}
\end{qbox}

\begin{solbox}
\textbf{Atomic Commit Protocol.}
An atomic commit protocol ensures that a distributed transaction that spans multiple
servers either \textbf{commits at all servers} or \textbf{aborts at all servers}.
No partial commit is allowed (= the Atomicity property of ACID in a distributed setting).
This is an instance of the \textbf{consensus problem}: all participating servers must
agree on one decision (commit or abort) despite possible message losses and crashes.

\medskip
\textbf{Two-Phase Commit (2PC).}

\textbf{Phase 1 — Prepare/Voting:}
\begin{enumerate}
  \item The \textbf{coordinator} sends a \texttt{PREPARE} message to all participating servers.
  \item Each server:
        \begin{enumerate}
          \item Writes its tentative updates to \textbf{stable storage} (disk) — before voting.
          \item Sends \texttt{YES} (ready to commit) or \texttt{NO} (cannot commit, e.g.,
                data corrupted or constraint violated) back to the coordinator.
        \end{enumerate}
\end{enumerate}

\textbf{Phase 2 — Decision:}
\begin{enumerate}
  \item \textbf{All YES received within timeout:}
        coordinator sends \texttt{COMMIT}; all servers apply updates from disk to permanent
        store and reply \texttt{OK}.
  \item \textbf{Any NO or timeout:}
        coordinator sends \texttt{ABORT}; all servers discard their tentative updates.
\end{enumerate}

\medskip
\textbf{2PC for Nested Transactions.}

A nested (or sub-) transaction is a transaction launched within the scope of a parent
transaction. The key property: a sub-transaction may commit locally, but its commit
is \emph{tentative} — if the parent transaction aborts, all sub-transactions are also
rolled back.

Protocol:
\begin{enumerate}
  \item \textbf{Phase 1 — Prepare (top-down):}
        The top-level coordinator sends \texttt{PREPARE} to each participating server.
        Each server that has running sub-transactions must first \emph{complete all
        sub-transactions locally}: committed sub-transactions are held as tentative;
        any failed sub-transaction triggers a local \texttt{NO} vote.
        Each server then votes \texttt{YES} (all sub-transactions succeeded and data is
        on stable storage) or \texttt{NO} (at least one sub-transaction failed).

  \item \textbf{Phase 2 — Commit or Abort (top-down):}
        \begin{itemize}
          \item If all servers vote \texttt{YES}: coordinator broadcasts \texttt{COMMIT};
                all servers make their tentative sub-transaction results permanent.
          \item If any server votes \texttt{NO}: coordinator broadcasts \texttt{ABORT};
                all servers roll back \emph{all} sub-transactions, including those that had
                committed locally.
        \end{itemize}
\end{enumerate}

\textbf{Key rule:} A locally-committed sub-transaction is \emph{not} final until the
parent transaction commits. The outer 2PC provides the global atomicity guarantee across
all participating servers and their nested sub-transactions.
\end{solbox}

\begin{notesbox}
\textbf{Exam pattern:} 2024 Q2c (3 marks).
Breakdown: $\sim$0.5 marks for atomic commit definition, $\sim$1 mark for basic 2PC
(Phase 1 + Phase 2), $\sim$1.5 marks for the nested-transaction extension.
Key point the examiner wants: locally-committed sub-transactions can still be rolled
back if the parent aborts — that's what makes 2PC for nested transactions distinct.
\end{notesbox}

\newpage
%======================================================================
\section{2024 — Q4c (3 marks)}
%======================================================================

\begin{qbox}
\textbf{Q4(c):} Briefly explain the different types of ordering of multicast messages
in overlapping groups. \hfill \textit{[3 marks]}
\end{qbox}

\begin{solbox}
In a distributed system with \textbf{overlapping process groups} (a process may belong to
more than one group simultaneously), multicast messages must be delivered in a consistent
order to prevent replicas in multiple groups from diverging.

\medskip
\textbf{1. FIFO Ordering.}
Messages sent by the \emph{same sender} are delivered at all receivers in the order
they were sent.
\begin{itemize}
  \item If process $P$ sends $M_1$ then $M_2$, all group members deliver $M_1$ before $M_2$.
  \item Messages from \emph{different senders} may arrive in any order.
  \item Problem for overlapping groups: process $B$ in both $G_1$ and $G_2$ sees
        messages from each group in sender order, but cannot guarantee global ordering
        across different senders.
\end{itemize}

\textbf{2. Causal Ordering.}
If message $M_1$ \emph{causally precedes} $M_2$ (i.e., $M_2$ was sent after $M_1$ was
received, or they share a causal chain), then all receivers deliver $M_1$ before $M_2$.
\begin{itemize}
  \item Stronger than FIFO: captures happen-before relationships across senders.
  \item Implemented using vector clocks.
  \item Still allows concurrent messages (not causally related) to be delivered in
        different orders at different nodes.
\end{itemize}

\textbf{3. Total Ordering (Atomic Broadcast).}
All processes in the group deliver \emph{all} messages in \textbf{exactly the same order},
regardless of sender or causal relationship.
\begin{itemize}
  \item Achieved via a sequencer (central token), consensus (Paxos/Raft), or
        a group ordering protocol.
  \item Critical for overlapping groups: a process $B$ in two groups $G_1$ and $G_2$
        will apply all updates in the same sequence as every other member of both groups.
  \item Required for active replication with Replicated State Machines — all replicas
        receive the same sequence of updates $\Rightarrow$ arrive at the same state.
\end{itemize}

\textbf{4. Hybrid Ordering (e.g., FIFO-Total or Causal-Total).}
Combines total ordering with FIFO or causal ordering.
Ensures both sender-order preservation \emph{and} global agreement on delivery sequence.
Used in active replication systems that also need causal consistency.

\medskip
\textbf{Why ordering matters for overlapping groups:}
Consider $G_1 = \{A, B\}$ and $G_2 = \{B, C\}$. Node $B$ belongs to both.
Without total ordering, $B$ may deliver a write to object $X$ from $G_1$ after $C$ has
already processed a read from $G_2$ that should have seen the update, causing
inconsistency. Total ordering ensures all nodes — including those in multiple overlapping
groups — agree on the same global delivery sequence.
\end{solbox}

\begin{notesbox}
\textbf{Exam pattern:} 2024 Q4c (3 marks).
The examiner wants: (1) FIFO — same sender, send order; (2) Causal — happen-before
preserved; (3) Total — same order at all receivers; and the overlapping-groups motivation.
3 marks → roughly 1 mark per ordering type. Keep definitions one sentence each, then
explain the relevance to overlapping groups.
\end{notesbox}

\newpage
%======================================================================
\section{Summary Table}
%======================================================================

\begin{center}
\renewcommand{\arraystretch}{1.35}
\begin{tabular}{lllll}
\toprule
\textbf{Year} & \textbf{Q} & \textbf{Marks} & \textbf{Topic} & \textbf{Key answer} \\
\midrule
2022 & Q3a & 3 & ACID crash: direct-to-disk & Slow I/O + partial write + no rollback \\
2022 & Q3b & 3 & ACID crash: batch-50-memory & Durability violated; crash loses 50 txns \\
2022 & Q3c & 3 & ACID crash: memory + log & Log grows; slow recovery; double-write \\
2024 & Q2a & 1.5 & Consensus + logical/physical & Agree on one value; logical=interleaved, physical=parallel \\
2024 & Q2c & 3 & Atomic commit + 2PC nested & Phase 1: prepare+vote; Phase 2: commit/abort; nested=local tentative \\
2024 & Q4c & 3 & Multicast ordering types & FIFO / Causal / Total / Hybrid; overlapping-group problem \\
\midrule
\multicolumn{3}{l}{\textbf{Total covered this PDF}} & \multicolumn{2}{l}{\textbf{16.5 marks}} \\
\bottomrule
\end{tabular}
\end{center}

\begin{notesbox}
\textbf{Pending (answered in L9 / L10 solutions):}
\begin{itemize}
  \item \textbf{Consistency models} (weak/causal/release/strong, quorums, linearizability):
        2020 Q4a/b, 2021 Q2a–c, 2022 Q4d, 2024 Q2b
        $\Rightarrow$ Covered in DSM (Lecture-09) wiki page and solutions.
  \item \textbf{ACID definition + deadlock + locking + WFG + timestamp ordering}:
        2021 Q1d, 2021 Q6a/b, 2022 Q6a, 2022 Q7a/b/c, 2024 Q3c, 2024 Q6a, 2024 Q8a/b
        $\Rightarrow$ Covered in RPC \& Concurrency (Lecture-10) solutions.
\end{itemize}

\textbf{Key terms for this lecture in exam:}
replication transparency, one-copy serializability, passive/active replication,
FIFO/Causal/Total ordering, virtual synchrony, atomic commit, Two-Phase Commit,
coordinator, prepare, voting, one-phase-commit shortcomings, write-ahead log.
\end{notesbox}

\end{document}
